\section{Faithfulness Scoring Is Metric-Fragile}
\label{sec:multimetric}

The previous two sections held the faithfulness scorer fixed at
DeBERTa-v3 NLI. Different scorers can disagree substantially on
the same answers; we re-score the entire $7{,}500$-row
$n{=}500{\times}5$-seed cell from \S\ref{sec:scaling} with two
additional scorers in parallel: roberta-large-mnli (a different
NLI model with independent training data) and a Mistral-7B-as-judge
RAGAS-style rubric.

\paragraph{Per-condition means across the three scorers.}

\begin{table}[H]
\centering
\small
\caption{Per-condition means across three faithfulness scorers on
the same $n{=}7{,}500$ rows. The paradox magnitude
$\overline{\phi}_{\text{baseline}} - \overline{\phi}_{\text{v1}}$
is reported in the right column.}
\label{tab:multimetric-means}
\begin{tabular}{lccccc}
\toprule
Condition & $n$ & DeBERTa-v3 & mnli & RAGAS judge & Magnitude (b-v1) \\
\midrule
Baseline   & 2500 & 0.6609 & 0.3501 & 0.7296 & --- \\
HCPC-v1    & 2500 & 0.6503 & 0.3184 & 0.5904 & --- \\
HCPC-v2    & 2500 & 0.6612 & 0.3509 & 0.7279 & --- \\
\midrule
\multicolumn{2}{l}{\textbf{Paradox magnitude:}} & $0.011$ & $0.032$ & $\mathbf{0.140}$ & \\
\bottomrule
\end{tabular}
\end{table}

The paradox magnitude is metric-dependent by an order of magnitude.
Under DeBERTa-v3 the magnitude is $0.011$; under
roberta-large-mnli, $0.032$; under the RAGAS-style LLM-judge,
$0.140$.

\paragraph{Pairwise agreement.}
Pairwise Pearson and Spearman correlations across all $n{=}7{,}500$
paired observations:

\begin{table}[H]
\centering
\small
\caption{Pairwise correlations among the three faithfulness
scorers on $n{=}7{,}500$ paired observations.}
\label{tab:multimetric-corr}
\begin{tabular}{llcc}
\toprule
Metric A & Metric B & Pearson $r$ & Spearman $\rho$ \\
\midrule
DeBERTa-v3 & roberta-large-mnli & $0.259$ & $0.265$ \\
DeBERTa-v3 & RAGAS judge        & $\mathbf{0.182}$ & $0.212$ \\
roberta-large-mnli & RAGAS judge & $0.674$ & $0.651$ \\
\bottomrule
\end{tabular}
\end{table}

DeBERTa-v3 and the RAGAS judge agree at Pearson $r=0.18$ on
faithfulness over the same $7{,}500$ paired observations. By
contrast, mnli and the RAGAS judge agree at $r=0.67$. DeBERTa-v3
is the outlier: under it the paradox magnitude is small
($0.011$); under either alternative the magnitude is several times
larger.

\paragraph{Two interpretations, both honest.}
\textbf{(a) The paradox is real and DeBERTa-v3 systematically
underestimates it.} Under the RAGAS judge, the paradox magnitude
($0.140$) is large, robust across seeds, and the recovery effect
($v2 - v1 = 0.138$) is correspondingly large. If the RAGAS judge is
the more accurate measurement, the published paradox claim is
supported but the published \emph{magnitude} is wrong because
DeBERTa-v3 understates it.

\textbf{(b) The paradox is small and the RAGAS judge inflates it.}
Under DeBERTa-v3 the paradox magnitude is $0.011$; under
roberta-large-mnli it is $0.032$. If those two NLI models are
better-calibrated for sentence-level entailment, the RAGAS-style
judge --- a larger generative model graded by a different
distribution --- may be picking up signals other than entailment
(answer fluency, surface plausibility, format).

\paragraph{Self-bias risk.}
The RAGAS judge in our setup uses a separate Mistral-7B instance,
which is the same model family as the system under test. We mitigate
direct identity by running the judge as an independent process with
a different prompt template, no shared context, and a structured
JSON-schema response, but the residual self-bias risk is real and
we report it as a limitation. A future version of this audit using
a frontier-scale judge (Claude-3.5-Sonnet, GPT-4o) would resolve
this --- but doing so honestly costs paid API tokens that the
zero-dollar constraint of this work disallows.

\paragraph{What this means for the literature.}
Across recent RAG-faithfulness papers, single-scorer evaluation
is the norm. Our finding has direct implications: any RAG
faithfulness claim with a magnitude under $\sim 0.05$ on a single
scorer should be reported alongside a second scorer with low
overall agreement to the first, and the magnitude should be
reported per-scorer. Not doing so risks publishing magnitudes
that are scorer-dependent without acknowledging it.
A 99-item two-rater human-eval template is staged at
\texttt{data/revision/fix\_03/human\_eval\_template.jsonl} for
follow-up work that wants to anchor the choice between
interpretations (a) and (b) in human judgement.
